
\section{Introduction}\label{sec:intro}

A sensory neuron turns a stimulus into a firing rate, and the cell it drives sees only how many spikes arrive in a short window. We call
the map from stimulus to expected spike count the tuning curve; for sound level it is the rate-level function. If the neuron is to
carry as much information as it can about the stimulus, which tuning curve should it have? Barlow asked this question in 1961 with the
stimulus distribution held fixed: the neuron cannot choose what the world presents, only how it responds~\cite{Barlow1961}. The tuning
curve that maximizes the mutual information between stimulus and count is called the infomax tuning curve. Laughlin gave the classic
answer for a neuron whose response is limited by its range but not by noise: the response should follow the cumulative distribution of
the stimulus, so that every response level is used equally often~\cite{Laughlin1981}. Laughlin's curve is smooth; for a neuron
whose count in a short window is noisy, we find that the answer is a staircase.

Spikes are noisy, and a short window holds few of them. The standard model of this noise is the Poisson count, whose variance
equals its mean; we build the theory on it and then test, on recorded fibres, how much their more regular counts change its
predictions. We consider a neuron whose count in the window is Poisson, with a mean that
lies between its spontaneous count $\lambda$, the count with no stimulus, and $\lambda+A$. The peak count $A$ is the largest
expected count above the spontaneous one, the peak firing rate times the window. Seen as a communication channel from rate to count,
such a neuron has a best distribution of rates, the one under which the count carries the most information. Rates closer together
than the width of the noise are hard to tell apart from the count, so the range is best used by a few well separated rates.
Information theory calls this best distribution the capacity-achieving input, and for this channel it indeed puts all its probability
on a few rates~\cite{Shamai1990}. Numerical work found that the infomax tuning curve of a
Poisson neuron is a staircase whose number of steps grows as the window lengthens~\cite{Nikitin2009}; in the opposite limit of many
spikes, the optimum approaches a smooth curve set by the stimulus distribution and the noise~\cite{McDonnellStocks2008}. Two questions
stayed open. How many steps does the optimal staircase have, as a function of the peak count? And do real neurons follow it?

This paper answers the first question with a theorem and puts the second to registered tests of two predictions: the shape of the
tuning curve's envelope, and the number of distinct rates the curve uses. We show that, for any continuous
stimulus distribution, the infomax tuning curve assigns the rates of the capacity-achieving input to the stimulus by its quantiles
(Proposition~\ref{prop:L1}). Quantiles are the stimulus values that cut its distribution into pieces of given probability, as the
median cuts it into halves. The stimulus distribution therefore decides where the steps sit and nothing else: their number and their
values depend only on $A$ and $\lambda$. We prove that the number of steps grows at least like $0.056\,A^{2/3}$
(Theorem~\ref{thm:C}). A naive count of resolvable rates, the range of rates divided by the width of the noise, grows only like
$\sqrt A$. Because the variance of a Poisson count equals its mean, the square root makes the noise nearly the same at every rate. In
the coordinate $2\sqrt{r}$ the Poisson noise has unit width, and the range from $0$ to $A$ becomes an interval of length $2\sqrt A$,
the Fisher length of the range, named after the Fisher information, which measures how precisely a count determines its rate. The optimal staircase uses more steps than that. The theorem concerns large peak counts; the fibres
tested below have peak counts of about 3 to 14 (5 to 14 for those of the shape test), where the optimal staircase for Poisson
counts has two to five steps. For large peak
counts the square root of the optimal rate follows the distribution function of the stimulus, which is Laughlin's rule with the
Poisson noise taken into account. The same problem with additive
Gaussian noise of fixed variance, the Gaussian channel, has a law for its number of levels, about $0.4\,a^{4/3}$ for an input range
of half-width $a$ in units of the noise, proved from below and measured~\cite{Paper1}. Carried over through the Fisher length, it gives
a sharper law for the exact number of steps: a new step appears each time the peak count passes a threshold, and the law predicts
these thresholds. The thresholds tested lie at peak counts of about 300 to 2000, far above those of neurons, so they test the
mathematics of the law and not its use for neurons. We state it as a conjecture, together with the
predictions it made before they were computed, including one that failed.

For the second question we used single auditory-nerve fibres of the gerbil, recorded at many sound levels with repeated
presentations~\cite{Heeringa2024}. The auditory nerve is a natural place to look. Its rate-level functions shift towards the most
frequent levels when the level distribution changes, in the midbrain and, more weakly, in the nerve itself~\cite{Wen2009}. The
decibel amplitude of natural sounds has a roughly symmetric bell-shaped distribution~\cite{Escabi2003}, for which the theory predicts a
specific asymmetric rate-level function. Every test was recorded, with its decision rule, in a dated log before the data were analysed.
The log is in the repository named under Data availability, and the entries for these tests carry a certificate of their date
from an independent authority, so the order of prediction and analysis can be checked (Section~\ref{sec:methods}). The failures are
reported with the successes. Pooled over fibres, the rate-level functions favour the predicted asymmetric shape over a symmetric sigmoid, in young and in aged animals.
Under Poisson counts about half of the fibres use more distinct rates, that is significantly different mean counts across sound
levels, than the optimal staircase allows, most by one. Their counts, however, are far less variable than
Poisson counts. For counts matched to the mean regularity of each fibre, more than four in five fibres stay within the optimal
staircase. For counts whose regularity is matched at every rate, the model that fits them best, about three in four do, short of the
registered $80\%$. A more powerful likelihood test finds about half of the fibres consistent with the optimal number of steps and nearly nine in ten
with one step more.

Kostal and Shinomoto~\cite{KostalShinomoto2016} studied a Poisson neuron with the same kind of bounds on its rate. Their downstream
reader sees the full spike train, and the best input then switches between two rates. Here the reader sees only the count in a
window, and counting is what makes the number of optimal rates grow with the peak count instead of staying at two.

Section~\ref{sec:results} holds the results. Section~\ref{sec:laughlin} derives the infomax tuning curve for a given stimulus
distribution. Section~\ref{sec:thmC} states the lower bound on the number of steps, whose proof is in Appendix~\ref{app:C}.
Section~\ref{sec:countlaw} states the conjectured law for the exact number and the predictions that tested it, and Section~\ref{sec:an}
reports the registered tests on auditory-nerve fibres. Section~\ref{sec:discussion} discusses them, and Section~\ref{sec:methods}
describes the data, the fits and the tests. The Supplementary Notes (S1 Text) treat counts more or less regular than Poisson counts, a
spontaneous rate, codes for windows that hold few spikes, and a test for discrete levels on six further public datasets.
